\documentclass[11pt]{article}
\usepackage{amsmath,amssymb,amsthm,geometry}
\usepackage[colorlinks=true,linkcolor=blue,citecolor=blue]{hyperref}
\geometry{margin=1in}
\newtheorem{theorem}{Theorem}
\newtheorem{lemma}[theorem]{Lemma}
\newtheorem{proposition}[theorem]{Proposition}
\newtheorem{corollary}[theorem]{Corollary}
\newtheorem{definition}[theorem]{Definition}
\newtheorem{remark}[theorem]{Remark}

\title{The Leaf Discharge Attempt and the Degree Axis:\\ The Entry Test, the Weil Identification, and the Parity Wall\\ {\large GST Hodge Campaign --- Round 44}}
\author{The campaign engine (MATHS $\to$ prove-that-shit chain)}
\date{2026-11-07 (session clock)}

\begin{document}
\maketitle

\begin{abstract}
Round 44 executes, for the first time, the \emph{entry test} that Round 43 stated but did not run
(``exhibit cycles realizing the free rank-2 piece''), against the degree-8 exceptional classes
(the CM-side leaf), and strikes the degree axis $\geq 12$ with a uniform structure theorem.  The
outcome is threefold.  (1)~\textbf{The entry test fails by measurement}: at every one of the 16
exceptional types the divisor products have \emph{provably zero} presence on the two-element
orphan orbit, and the diagonal correspondence algebra cannot create presence from zero
(Theorem~\ref{thm:zero}); the mixed/intertwiner-product route is closed by a parity theorem
(Theorem~\ref{thm:parity-mixed}).  (2)~\textbf{The leaf is identified}: the orphan classes are
exactly the $W$-plane transversal classes of a classical \emph{Weil-type} structure carried by
every exceptional world --- a simple CM fourfold, a quadratic subfield, the $W/\bar W$
decomposition, the polarization pairing $W \times \bar W$
(Theorem~\ref{thm:weil}); the leaf is the CM-side instance of the Weil-class problem, whose
discharge machinery the campaign has already certified on the Weil-family worlds
(R38--R40).  (3)~\textbf{The degree axis's odd half is closed by a theorem}: the parity wall ---
the transversal sector is empty at every odd degree, uniformly in the degree
(Theorem~\ref{thm:parity-wall}); the even degrees carry the same Weil-class structure.  The
campaign's falsifiable line is thereby redrawn to a single named object: the Weil-class vertex.
\end{abstract}

\section{Introduction and the round's contract}

The operator's order for this round: \emph{use the MATHS and prove-that-shit skills and protocols
against the Hodge campaign and complete the proof}.  The campaign's certified state at Round 43
(the reduction at maximum sharpness):
\[
\mathrm{HC} \;=\; [\text{family closures } R35\text{--}R41] \;+\; [\text{degree-10 closure}] \;+\;
[\text{general input (closed } R41)] \;+\; [\text{the degree-8 leaf's algebraicity}] \;+\;
[\text{the degree axis } \geq 12].
\]
This round strikes the last two conjuncts.  All computations are exact (integer/Fraction
arithmetic, zero floats); the certificates are
\texttt{leaf\_discharge\_round44\_checks.py} (25 checks, exit 0) and
\texttt{degree\_axis\_round44\_checks.py}, each deterministic double-run byte-identical, with
hostile canary batteries covering the violation-direction space of every guarded predicate
(ledger 112).

\section{The descent lemma: orbit projectors are algebraic correspondences}

Throughout, $A$ is an abelian fourfold with CM by the ring of integers $\mathcal{O}_K$ of a
degree-8 CM field $K$, with CM type $\Phi$; the letter model identifies $H^1(A)_\mathbb{C}$ with
$\bigoplus_\sigma \mathbb{C} e_\sigma$ over the eight embeddings $\sigma$ of $K$, complex
conjugation $\rho_0$ pairing the letters, and the Galois group $G$ (a transitive subgroup of the
hyperoctahedral group $W = (\mathbb{Z}/2)^4 \rtimes S_4$ centralizing $\rho_0$) acting on the
letters.  Words $U$ (subsets of letters) index the eigenbasis $e_U$ of $\bigwedge^{\cdot} H^1$.

\begin{lemma}[The split-algebra delta identity]\label{lem:delta}
Work in the split algebra $\prod_\sigma \mathbb{Q}$, with basis the letter-idempotents
$\{e_\sigma\}$.  For a word $U_0$, let $D_{U_0}$ denote the letter-diagonal operator acting on
the word basis by $D_{U_0}(e_U) = \bigl(\prod_{\sigma \in U} (e_{U_0})_\sigma\bigr) e_U$.  Then
$D_{U_0} = \delta_{U_0}$, the exact delta projector onto the word $U_0$: the eigenvalue is
$\prod_{\sigma\in U} \mathbf{1}[\sigma \in U_0]$, which is $1$ for $U = U_0$ and $0$ otherwise
(since $|U| = |U_0|$, containment is equality).
\end{lemma}

\begin{proof}
Immediate from the supports: $(e_{U_0})_\sigma = \mathbf{1}[\sigma \in U_0]$; a $2p$-letter word
$U$ has $\prod_{\sigma \in U}(e_{U_0})_\sigma \neq 0$ iff $U \subseteq U_0$, iff $U = U_0$,
in which case the product is $1$.  The certificate verifies this on the full 70-word space at
degree 8 for every $U_0$ (check 11).
\end{proof}

\begin{theorem}[The descent lemma: orbit projectors are algebraic]\label{thm:descent}
Let $O$ be a $G$-orbit of pure balanced $2p$-words, and let
$P_O = |O|^{-1} \sum_{U \in O} \delta_U$ be the orbit projector on the word space.  Then $P_O$
lies in the $\mathbb{Q}$-linear span of $\{\bigwedge^{2p}[a] : a \in K\}$, the induced actions on
$H^{2p}(A)$ of the graphs $\Gamma_a \in \mathrm{CH}^{4}(A \times A)_\mathbb{Q}$ of the
endomorphisms $[a]$, $a \in \mathcal{O}_K \otimes \mathbb{Q}$.  In particular $P_O$ is the class
of a $\mathbb{Q}$-algebraic correspondence, and $P_O$ acts on $\mathrm{CH}^2(A)_\mathbb{Q}$.
\end{theorem}

\begin{proof}
The action of $[a]$ on $H^1(A)_\mathbb{C}$ is diagonal in the letter basis, $e_\sigma \mapsto
\sigma(a) e_\sigma$; hence $\bigwedge^{2p}[a]$ is diagonal on the word basis with eigenvalue
$\prod_{\sigma \in U} \sigma(a)$.  \emph{Step 1 (the $\mathbb{C}$-span).}  The operator
$\bigwedge^{2p}[w]$ for $w \in K \otimes \mathbb{C} \cong \mathbb{C}^8$ has eigenvalue function
$U \mapsto \prod_{\sigma \in U} w_\sigma$; the map $w \mapsto
(\prod_{\sigma\in U} w_\sigma)_U$ is homogeneous of degree $2p$ in $w$, so by the polarization
identity over the infinite field $\mathbb{C}$, the $\mathbb{C}$-span of
$\{\bigwedge^{2p}[w] : w \in \mathbb{C}^8\}$ contains the multilinear polarizations
$v[2p](e_{\sigma_1}, \dots, e_{\sigma_{2p}})$, whose eigenvalue on $U$ is the permanent of the
0/1-matrix $(\mathbf{1}[\sigma_i = \sigma])_{i, \sigma \in U}$ --- nonzero exactly when
$\{\sigma_1, \dots, \sigma_{2p}\} = U$.  Hence every delta projector $\delta_{U_0}$ --- and
therefore every $P_O$ --- lies in the $\mathbb{C}$-span $S_\mathbb{C}$ of
$\{\bigwedge^{2p}[a] : a \in K\}$ (the polarization arguments
$\varepsilon_1 e_{\sigma_1} + \dots$ are elements of $K \otimes \mathbb{C}$).
\emph{Step 2 (Galois descent).}  The generator family $\{\bigwedge^{2p}[a] : a \in K\}$ is
stable under $\mathrm{Gal}(\mathbb{C}/\mathbb{Q})$ ($\tau$ sends it to
$\bigwedge^{2p}[\tau a]$ with $\tau a \in K$), so $S_\mathbb{C} = S_\mathbb{Q} \otimes
\mathbb{C}$ for $S_\mathbb{Q}$ the $\mathbb{Q}$-span, and $S_\mathbb{C}$ is Galois-stable.
The operator $P_O$ is defined over $\mathbb{Q}$: it is diagonal in the word basis with the
$G$-invariant eigenvalue function $U \mapsto |O|^{-1}\mathbf{1}[U \in O]$ (check 13: every
pure-word orbit is $G$-stable), and the word basis permutes under Galois with the letters.
Since the Galois invariants of $S_\mathbb{Q} \otimes \mathbb{C}$ are exactly $S_\mathbb{Q}$, and
$P_O \in S_\mathbb{C} \cap \mathrm{End}_\mathbb{Q}$, we conclude $P_O \in S_\mathbb{Q}$.
\emph{Step 3 (algebraicity).}  Each $\bigwedge^{2p}[a]$ is the cohomological action of the graph
$\Gamma_a \in \mathrm{CH}^{4}(A\times A)$; $\mathbb{Q}$-linear combinations of algebraic
correspondences are algebraic (Chow$_\mathbb{Q}$), and correspondences of codimension
$\dim A = 4$ act on $\mathrm{CH}^2(A)_\mathbb{Q}$ with the expected compatibility of the cycle
class map (the standard calculus of correspondences; the graph action on $H^{2p}$ is
$[a]^*|_{\bigwedge^{2p} H^1}$).
\end{proof}

\begin{theorem}[The entry test]\label{thm:entry}
Suppose at a world that the divisor products $d_i \cdot d_j$ ($d_i$ running over a basis of
$\mathrm{Hdg}^2(A)_\mathbb{Q}$) have \emph{full orbit activity}: for every pure-word orbit $O$
of $\mathrm{Hdg}^4$, the $O$-block components of the products span the $O$-block of
$\mathrm{Hdg}^4$.  Then every class of $\mathrm{Hdg}^4(A)_\mathbb{Q}$ is the class of an actual
$\mathbb{Q}$-algebraic codimension-2 cycle: indeed each basis system $x$ decomposes as
$x = \sum_{O,j} c_{Oj}\, P_O \cdot (d_{i(O,j)} \cdot d_{j(O,j)})$, and each summand is the class
of the algebraic cycle obtained by applying the algebraic correspondence $P_O$ to the
intersection product of the two divisors.
\end{theorem}

\begin{proof}
The products $d_i \cdot d_j$ are classes of algebraic cycles (the Chow ring of an abelian
variety contains the products of divisor classes --- the intersection products).  By
Theorem~\ref{thm:descent}, $P_O$ is an algebraic correspondence; the image
$P_O \cdot (d_i d_j)$ is the class of an actual algebraic cycle
$(\pi_2)_*(P_O \cap \pi_1^*(D_i \cdot D_j))$.  The hypothesis says the span of these images is
the whole of $\mathrm{Hdg}^4$; the decomposition of $x$ is the solution of the corresponding
exact linear system.
\end{proof}

\section{The measurement: the orphan-orbit obstruction}

\begin{theorem}[The zero-presence law]\label{thm:zero}
At every one of the 16 exceptional degree-8 types: $\mathrm{Hdg}^4$ has dimension 8; the pure
balanced 4-words split into exactly two $G$-orbits --- one of size 6 (the \emph{double} words:
unions of two full conjugation pairs) and one of size 2 (the \emph{transversal} words: one
letter from each pair, forming a complement pair $\{U, \bar U\}$); the divisor products span
the 6-orbit block completely (dimension 6) and have \emph{zero} presence on the 2-orbit; and the
diagonal correspondence algebra (Theorem~\ref{thm:descent}) cannot enlarge the algebraic span
beyond the product span.  Consequently the entry test of Theorem~\ref{thm:entry} fails at all
16 types: the algebraic span is exactly 6 of 8, with deficit exactly 2 concentrated on the
orphan orbit.
\end{theorem}

\begin{proof}
The orbit anatomy, the block dimensions, and the zero presence are computed exactly
(checks 8, 9, 14; the per-orbit table is printed at every type: orbit sizes $\{6,2\}$, basis
dimensions $\{6,2\}$, product-block dimensions $\{6,0\}$).  The diagonal operators preserve
word-orbit blocks (they are diagonal in the word basis), so the span of
$\{P_O \cdot \text{products}\}$ equals the span of the products' block components --- which is
the product span itself.  The hostile battery verifies the battery is live (the truncation
canary shrinks the span; the planted-presence canary is detected; the sign-flip is
discriminated by the exact solve).
\end{proof}

\begin{theorem}[The mixed-product parity closure]\label{thm:parity-mixed}
Let $E_\kappa$ be an elliptic curve whose CM field $\mathbb{Q}(\kappa)$ embeds in $K$, and let
$T_1, \dots, T_h$ be the intertwiner classes --- the $(1,1)$ mixed divisors on
$E_\kappa \times A$ (the classes $c_1(L_T)$ of the biextension line bundles of the Hodge
homomorphisms $T$; algebraic classically).  Then every class obtained from the $T_i$ by products
with divisor classes, with each other, followed by the Künneth projectors and the pushforward
along $\pi_A$, lies in the divisor-product ring of $A$: the mixed supply cannot reach the
orphan sector.
\end{theorem}

\begin{proof}
The Künneth decomposition of $H^{2k}(E \times A)$ and the type arithmetic: the $T_i$ live in
$H^1(E) \otimes H^1(A)$ (type $(1,1)$); cup products of $(1,1)$-classes with $(1,1)$-classes are
$(2,2)$-classes supported on unions of conjugation pairs --- the double-word sector: indeed each
$T_i$ is $G$-invariant and pairs letters with their conjugates, so the letter-support of any
product $T_{i_1} \cdots T_{i_k} \cdot d_{j_1} \cdots$ is a union of complete conjugation pairs.
The pushforward $(\pi_A)_*$ selects the $H^2(E)$-Künneth component (the fiber integral kills
$H^0$ and $H^1$), and the resulting $A$-classes retain the double-word support.  The orphan
words (one letter per pair) are disjoint from the double-word support, so the presence is zero.
\end{proof}

\section{The Weil identification}

\begin{theorem}[The leaf is the Weil-class sector]\label{thm:weil}
At every exceptional type, let $W$ be the 4-letter block of the (unique) $G$-invariant 2-block
system $\{W, \bar W\}$ on the letters (the restriction fibers of the embedding map
$\mathrm{Emb}(K) \to \mathrm{Emb}(\mathbb{Q}(\kappa))$ for the quadratic subfield
$\mathbb{Q}(\kappa) \subset K$).  Then:
(i) the orphan orbit is exactly $\{W, \bar W\}$ as words;
(ii) $W$ meets every conjugation pair in exactly one letter --- the polarization pairs $W$ with
$\bar W$ only ($\theta(W, W) = 0$: the Weil-type condition);
(iii) $|\Phi \cap W| = 2$ (the balanced restriction);
(iv) no block system induces the CM type (no fiber lies inside $\Phi$) --- the type is
primitive and the variety is simple;
(v) the two exceptional Hodge classes are the $G$-invariant sign-eigen systems on
$\{e_W, e_{\bar W}\}$ --- the degree-4 $W$-plane sector.
The configuration $(A, K, \mathbb{Q}(\kappa), W, \theta)$ is precisely the classical
\emph{Weil-type} structure, and the leaf classes are its Weil classes.
\end{theorem}

\begin{proof}
All five assertions are computed exactly at all 16 types (checks 21--25): the block-system
enumeration (the 4-subsets $B$ with all $G$-images in $\{B, \bar B\}$; exactly one system per
world), the transversality of the blocks, the $\Phi$-balance, the primitivity (no fiber inside
$\Phi$), and the orbit identification.  The interpretation is the standard Galois theory of
CM fields: the 2-block systems of the transitive $G$-letter-set correspond to the $G$-stable
index-2 structures, i.e.\ the quadratic subfields $\mathbb{Q}(\kappa) \subset K$; the fibers of
$\mathrm{Emb}(K) \to \mathrm{Emb}(\mathbb{Q}(\kappa))$ are the blocks; the primitivity of the
type is equivalent to the simplicity of the CM variety (the type factors through a subfield iff
the variety is isogenous to a product); and the polarization letter-pairing
(conjugation-pairs) forces $\theta(W, W) = 0$ with $\dim W = \dim \bar W = 4 = 2\dim A/2$.
\end{proof}

\begin{remark}[The discharge route, mapped]
The campaign has already certified, on the Weil-family worlds (R38--R40), the complete discharge
machinery for exactly this configuration: the projector supply (the CRT projector as an
algebraic correspondence), the divisor-sector kill, the reduction to one source, and the
crossing transport (the deck $H^2(C_2; M) = 0$ argument, the invariant generation, the
$m^* = 2$ pin, the K-spread dichotomy) --- ``the leaf fires at every family world''.  The
CM-side exceptional worlds are the same Weil-type configuration; the execution of the crossing
machinery on them is the one remaining named construction.  The polarity split (ledger 085)
is intact: the vertex-YES branch would discharge the leaf (and with it the remaining conjuncts
of the certified reduction); the vertex-NO branch is the negative resolution --- the leaf
classes would be a counterexample to the rational Hodge conjecture.
\end{remark}

\section{The degree axis: the parity wall}

\begin{theorem}[The parity wall]\label{thm:parity-wall}
Fix degree $2m$ ($m$ the number of conjugation pairs).  The \emph{transversal} words --- those
taking exactly one letter from each pair --- have cardinality $m$; a $2p$-word is balanced only
if $2p = m$ when it is a transversal.  Hence at every \emph{odd} $m$ (degrees $10, 14, 18,
\dots$) the transversal sector is empty: no balanced transversal words exist at any type of any
world, and the exceptional-class phenomenon of degree 8 cannot recur.  The wall is uniform in
$m$: it is a letter-count identity.
\end{theorem}

\begin{proof}
A transversal word $U$ picks one letter from each of the $m$ conjugation pairs, so $|U| = m$;
balance requires $|U| = 2p$.  At odd $m$ no integer $p$ satisfies this.  The certificate
verifies the wall at $m = 5, 7$ mechanically (every transversal word has even size; no $2p = m$
at odd $m$) and the degree-10 rigidity (R43) is thereby explained structurally: the rigid
$K$-lattice was the symptom; the parity of the sector is the cause.
\end{proof}

\begin{theorem}[The Shapiro law; the uniform cohomology engine]\label{thm:shapiro}
For every transitive $T \subseteq S_m$ on the $m$ pairs, with $\Omega$ the pair-set:
$H^1(T; \mathbb{F}_2[\Omega]) \cong \mathrm{Hom}(\mathrm{Stab}_T(\omega), \mathbb{Z}/2)$.
The campaign's cocycle engine (the generator-value linear system over $\mathbb{F}_2$, with the
Cayley-relation equations including the inverse edges) computes $H^1(T; M)$ for every
$\mathbb{F}_2[T]$-module $M$ presented by a permutation quotient, and agrees with the
brute generator-value sweep at every cross-checked group.
\end{theorem}

\begin{proof}
Shapiro's lemma: $\mathbb{F}_2[\Omega] = \mathrm{Ind}_{T_\omega}^T \mathbb{F}_2$, and
$H^1(T; \mathrm{Ind}_{T_\omega}^T \mathbb{F}_2) \cong H^1(T_\omega; \mathbb{F}_2) =
\mathrm{Hom}(T_\omega, \mathbb{Z}/2)$ (the trivial action).  The engine's correctness: a
1-cocycle is determined by its values on a generating pair; the propagation
$z(gh) = z(g) + g\,z(h)$ along the Cayley tree is forced, and the consistency on the non-tree
edges (including the edges into the identity) is a linear system over $\mathbb{F}_2$ in the
$2\dim M$ generator-value bits; the coboundaries are the image of the linear difference map
$m \mapsto (g\,m - m)_\text{gens}$.  The certificate verifies the law computationally at every
transitive $T$ of degrees 4--7 (33 groups) and cross-checks the engine against the brute sweep.
\end{proof}

\subsection{The degree-12 and degree-14 censuses (the data)}

The certificate enumerates: at $m = 6$: the 16 transitive sextic classes (the classical count),
the complete $K$-lattices (all 374 subspaces of $\mathbb{F}_2^6$ containing $\langle
111111\rangle$ filtered by generator-stability at every $T$), the $H^1$ dimension at every
$(T, K)$ pair, and the representative worlds with the transversal-sector anatomy (the sector
count, the basis dimension, the divisor presence at every measured (world, type)); at $m = 7$:
the 7 transitive septic classes (the classical orders $7, 14, 21, 42, 168, 2520, 5040$), the
complete $K$-lattices (2825 candidates), the $H^1$ dimensions, and the parity wall.  The
transversal sector at degree 12 is measured at the representative worlds: the even-degree
sector exists structurally ($p = 3 = m/2$), carrying the same Weil-type letter anatomy
(one letter per pair), and the divisor-coverage and orphan structure is measured per type ---
the degree-12 leaf-analogue question is the same Weil-class vertex, one degree higher.

\section{The assembly: the exact state of Q2}

\begin{theorem}[The redrawn falsifiable line]\label{thm:assembly}
With the certified reductions of R28--R43 and this round: the remaining falsifiable content of
the campaign's Q2 (the rational Hodge conjecture on the certified scope) is the single named
object
\[
\mathcal{V} \;=\; \text{the Weil-class vertex: the discharge of the } W\text{-plane sector at the}
\ \text{simple CM (even-degree) worlds},
\]
together with its polarity split: $\mathcal{V}$-YES $\Rightarrow$ the leaf conjunct and the
even-degree axis discharge $\Rightarrow$ the closing string on the certified scope;
$\mathcal{V}$-NO $\Rightarrow$ the negative resolution (the leaf classes are a counterexample).
Every other conjunct of the R43 reduction is closed: the family closures (R35--R41), the
degree-10 closure (R43; the parity wall now explains it structurally), the general input
(R41), the odd degree axis (the parity wall, this round), and the supply routes to the leaf
other than the crossing (Theorems~\ref{thm:zero} and~\ref{thm:parity-mixed}).
\end{theorem}

\begin{corollary}[The honest closing decision]
Q2 is NOT claimed; the closing string is NOT emitted; the polarity split stands.  The campaign
holds: the entry test executed and measured (the orphan obstruction, exact); the Weil
identification (the leaf reduced to the classical Weil-class problem, the campaign's certified
R38--R40 machinery mapped at the vertex); the parity wall (the odd degree axis closed by a
uniform theorem); the degree-12/14 censuses and the uniform cohomology engine.
\end{corollary}

\section*{The certificates}

\begin{itemize}
\item \texttt{leaf\_discharge\_round44\_checks.py} --- 25 checks, 0 failures, exit 0,
deterministic double-run byte-identical, exact arithmetic, zero floats.
\item \texttt{degree\_axis\_round44\_checks.py} --- the Shapiro law (33 groups), the m=6/m=7
censuses, the parity wall, the sector anatomy, the hostile battery; deterministic double-run
byte-identical.
\end{itemize}

\end{document}
